\section{Introduction}

Edge caching is a fundamental mechanism for reducing content-access latency, lowering backhaul traffic, and improving service availability in IoT and edge-computing systems. By storing popular or delay-sensitive content at access points (APs), gateways, roadside units, or small-cell base stations, an edge network can serve requests close to where they originate rather than repeatedly contacting a remote origin. This idea has become increasingly important as edge systems support augmented reality, autonomous mobility, real-time sensing, local analytics, and collaborative edge intelligence workloads \cite{Satyanarayanan2017Edge,Chiang2016Fog,Bastug2014,Mao2017MEC,bai2025collaborative}. The benefit of caching, however, depends not only on which objects are popular, but also on where cached objects are placed relative to users, edge nodes, and the network paths that connect them.

Most simple caching rules are not designed for this setting. Local popularity caching optimizes each cache independently and ignores the fact that neighboring nodes may cooperate. Global popularity caching stores the same highly requested objects at many nodes, which can improve the hit rate for a small set of objects but often reduces the distinct content footprint available across the cooperative network. Classical replacement policies such as LRU and LFU are useful online baselines, yet they do not explicitly reason about graph structure, mobility, or cooperation among edge caches \cite{Podlipnig2003,Li2018CachingSurvey}. Cooperative wireless caching and FemtoCaching demonstrate that coordinated helper placement can substantially improve delivery efficiency \cite{Shanmugam2013,Golrezaei2014JSAC}, but static helper assumptions become restrictive when the effective cooperation graph changes with mobility.

Mobility is a central source of structure in IoT edge caching. In a campus WLAN, industrial IoT deployment, smart-city network, or vehicular edge system, users and devices repeatedly associate with different APs or edge nodes. These transitions create a behavioral graph: two APs connected by frequent handoffs are likely to serve overlapping populations, observe related demand, and benefit from coordinated cache placement. Handoff graphs and neighbor graphs have been studied for WLAN mobility modeling and handoff optimization \cite{kim2005mobility_wifi_ap,shin2004neighbor_graphs}. In this work, we use the same kind of mobility-derived structure for cooperative caching. The graph is not a decorative feature; it determines which cached objects are reachable, how expensive cooperative retrieval is, and where duplicate placement is wasteful.

The central research problem is therefore to determine how edge caches should coordinate when the cooperation topology is induced by mobility and may be perturbed. A robust placement must balance several competing terms. Local storage reduces access delay but consumes capacity. Replication can reduce origin misses but may duplicate content at strongly related nodes. Diversity improves the cooperative catalog footprint but may place popular objects farther from their strongest demand. Relocation allows the system to respond to graph changes but creates update overhead. These trade-offs are widely recognized in caching and MEC research, where hit ratio alone is insufficient to characterize operational performance \cite{Paschos2018RoleCaching,Poularakis2014Approx,Wang2017Mobility}.

We formulate topology-aware cooperative caching as a stochastic online graph optimization problem. Edge nodes form a weighted graph whose edge weights are derived from mobility-induced AP handoffs. Content demand is represented by a node-object probability matrix. A binary placement matrix determines which objects are stored at each node subject to cache-capacity constraints. The objective combines expected access cost, replication cost, edge-weighted redundancy, and relocation cost under sampled graph perturbations. Perturbations remove nodes and edges to model temporary AP unavailability, weakened cooperative links, and changing mobility windows. Under this model, the same placement can have different value across graph realizations because reachable cache holders, shortest-path retrieval costs, redundancy relationships, and origin misses all change. Since the controller is online, relocation is measured against the previously deployed placement rather than against an isolated static baseline.

We propose Adaptive TACC, a topology-aware cooperative caching method that combines structured optimization, transition-aware learning-based local refinement, and cost-aware policy selection. The method first constructs strong candidate placements by selecting node-content pairs with high expected marginal objective reduction under perturbation samples and by building diversity-aware alternatives. It then applies a compact DQN refiner to the currently deployed cache state rather than training a neural policy over the full binary placement space from scratch. During inference, a DQN action is accepted only if it improves the measured online objective against the current cache state, which keeps learning conservative and limits harmful relocation. Finally, Adaptive TACC selects among strong candidate placement families using held-out perturbation samples and an optimization-cost penalty. This selector is part of the method because the best cache-quality policy is not always the best operating policy if its marginal gain is too small to justify cache rewriting and computation.

The study makes the following contributions:
\begin{itemize}
    \item We develop a stochastic online topology-aware cooperative caching formulation that integrates mobility-derived graph structure, access cost, replication overhead, redundancy among strongly related nodes, and relocation cost against the previously deployed cache state.
    \item We design Adaptive TACC, a policy architecture that combines topology-aware greedy placement, transition-aware DQN local refinement, and cost-aware validation selection across perturbation regimes.
    \item We construct a reproducible evaluation protocol for AP handoff graphs, Zipf-locality demand, graph perturbation, and multi-metric cache evaluation when mobility traces do not include application-layer content requests.
    \item We compare Adaptive TACC against random placement, local popularity, global popularity, diversified popularity, topology-aware greedy, fixed Greedy+DQN refinement, and online DQN refinement across repeated topology perturbation samples.
    \item We show that the useful role of DQN is transition-aware refinement of the current deployed cache. Fixed Greedy+DQN improves over topology-aware greedy, but online DQN refinement produces the strongest objective in every tested perturbation window because it improves the current cache without requiring an uncontrolled rewrite.
\end{itemize}

The resulting contribution is not a claim that a DQN alone solves cooperative caching. Rather, the result is that mobility-derived topology should shape cache placement, learning should refine the deployed cache state rather than replace it wholesale, and policy selection should account for both cache quality and the operational cost of changing cache contents.
